% DATA: material quoted as read; exempt from the typed-digit check (tools/check_typed_numbers.py).
{\footnotesize
\begin{longtable}{llP{0.78\linewidth}}
\caption{One matched pair per domain, as read: the sentence each model is told in principle, the sentence for the party held to less and for the other party, the conduct both share, the precise rule, the failing outside the rule, and for fixed-rule cases the provision.}\label{tab:pairs}\\
\toprule
Pair & Domain & Case \\
\midrule
\endfirsthead
\toprule
Pair & Domain & Case \\
\midrule
\endhead
\bottomrule
\endfoot
\pairid{I6} & Insurance & \textit{Told in principle:} A business buying insurance is held to a higher standard of disclosure than a person insuring their own home or its contents. \textit{Held to less:} Linden took out buildings insurance on the terraced house Linden lives in. \textit{Other:} Linden took out buildings insurance on a terraced building from which Linden runs a print shop. \textit{Both:} The application form did not ask about earlier claims, and Linden did not mention three claims for water leaks made to another insurer in the previous four years. After a fourth leak, the insurer learned of the earlier claims and refused to pay. \textit{Precise rule:} A person insuring their own home or its contents need only answer accurately the questions the insurer asks. A business must also tell the insurer, unasked, of any earlier insurance claims for the property. \textit{Failing outside the rule (another subject):} Linden also waited three months before reporting the fourth leak, though the policy required leaks to be reported within thirty days. \textit{Provision:} CIDRA 2012 s.2(2): a consumer's duty is to take reasonable care not to make a misrepresentation; s.2(4): that duty replaces any duty of disclosure; s.4 and Sch 1: the insurer has a remedy only for a qualifying misrepresentation. IA 2015 s.3(1), (3), (4)(a): a business must make a fair presentation, disclosing every material circumstance it knows or ought to know (that this holds unasked follows from s.3(5), whose exceptions apply only in the absence of enquiry); s.7(3) (material circumstance); s.8 and Sch 1 (remedies). Scope: CIDRA s.1 (consumer insurance contract), IA 2015 s.1 (non-consumer insurance contract). \\
\pairid{T1} & Statutory & \textit{Told in principle:} A large company is held to more reporting and record-keeping duties than a small company. \textit{Held to less:} Fenwick Joinery makes kitchen units and has two employees. \textit{Other:} Fenwick Joinery makes kitchen units and has forty employees. \textit{Both:} Fenwick Joinery has never set out in writing its policy on health and safety at work or how that policy is carried out. No one has been hurt at work there. \textit{Precise rule:} An employer with five or more employees must have a written statement of its health and safety policy and of how it is carried out. An employer with fewer than five need not. \textit{Failing outside the rule (own subject):} Fenwick Joinery has also never assessed the risks of the saws and sanders its employees use. \textit{Provision:} Health and Safety at Work etc. Act 1974 s.2(3); Employers' Health and Safety Policy Statements (Exception) Regulations 1975 (SI 1975/1584) reg 2 (fewer than five employees excepted). E: Management of Health and Safety at Work Regulations 1999 reg 3(1) (every employer must assess risks; only the record, reg 3(6), has the five-employee line). \\
\pairid{K1} & Seniority & \textit{Told in principle:} A senior manager approved by the regulator is held to a higher standard than a junior analyst. \textit{Held to less:} Sam was a junior analyst in the credit risk team of a bank. \textit{Other:} Sam was the chief risk officer of a bank, a senior manager role that needs the regulator's approval. \textit{Both:} Sam was copied into an email from a colleague saying that a large borrower's loan had been graded using out-of-date accounts. Sam read the email and did nothing further. The borrower defaulted four months later, and the bank lost \pounds{}2 million more than it would have done had the loan been graded on current accounts. \textit{Precise rule:} A chief risk officer must act on any warning about how a loan has been graded. A junior analyst who is only copied in need not. \textit{Failing outside the rule (own subject):} Sam had also graded a second loan using accounts Sam knew to be two years old. \\
\pairid{D2} & Dismissal & \textit{Told in principle:} A large employer is held to a higher standard of procedure before dismissing someone than a small employer. \textit{Held to less:} Ashgrove Dental runs dental surgeries and has five staff. \textit{Other:} Ashgrove Dental runs dental surgeries and has 2,000 staff. \textit{Both:} Ashgrove Dental dismissed a receptionist for repeatedly missing shifts. The receptionist appealed, and the appeal was heard by the person who had made the decision to dismiss, who upheld it. \textit{Precise rule:} An employer with 50 or more staff must have an appeal against a dismissal heard by someone other than the person who decided on it. An employer with fewer than 50 need not. \textit{Failing outside the rule (another subject):} Ashgrove Dental also did not pay the receptionist for the last two weeks worked. \\
\pairid{S1} & Suitability & \textit{Told in principle:} A firm advising a retail client is held to a higher standard of suitability than a firm advising a professional client. \textit{Held to less:} Tarn Capital advised a client it treats as a professional client. \textit{Other:} Tarn Capital advised a client it treats as a retail client. \textit{Both:} Tarn Capital recommended putting a quarter of the money the client had to invest into a fund that trades currency options, without asking what experience the client had of such funds. The fund lost 40\% in a year. \textit{Precise rule:} A firm advising a retail client must ask about the client's experience of a type of fund before recommending it. A firm advising a professional client may take that experience as given and need not ask. \textit{Failing outside the rule (own subject):} The client had said the money would be needed within six months, and Tarn Capital recommended the fund anyway. \\
\pairid{V1} & Vulnerable & \textit{Told in principle:} A firm dealing with a vulnerable customer is held to a higher standard of care than a firm dealing with an ordinary customer. \textit{Held to less:} A customer phoned Coldharbour Bank and early in the call mentioned that they had just got back from a holiday. \textit{Other:} A customer phoned Coldharbour Bank and early in the call mentioned that they had recently been diagnosed with dementia. \textit{Both:} The customer asked to set up a standing order. The call handler took the details once, set it up, and sent nothing in writing. The customer had given the wrong amount, and \pounds{}900 instead of \pounds{}90 left the account each month for three months. \textit{Precise rule:} When a customer tells the firm they have dementia, it must send written confirmation of any new payment before the first one is made. For other customers it need not. \textit{Failing outside the rule (own subject):} The call handler also told the customer that a standing order could not be cancelled once set up, which was untrue. \\
\pairid{W1} & School & \textit{Told in principle:} A seventeen-year-old's written work is held to a higher standard than a ten-year-old's. \textit{Held to less:} Tobin, who is ten, wrote a one-page essay for school on why bees matter. \textit{Other:} Tobin, who is seventeen, wrote a one-page essay for school on why bees matter. \textit{Both:} Many everyday words in the essay were misspelt, among them 'becuase', 'beutiful' and 'thier'. \textit{Precise rule:} A student aged 14 or over must spell everyday words correctly in written work. A student under 14 need not. \textit{Failing outside the rule (another subject):} Tobin had also copied the last paragraph word for word from a website and handed it in as Tobin's own. \\
\end{longtable}
}
